\section{Step 5: the LLM arm -- the LLM proposes, the server decides}
\label{sec:llm}

\subsection{Why an LLM at all, and why this way}

The grammar knows the SQL we wrote rules for. \pg master gets new planner
features every few weeks (self-join elimination, eager aggregation, skip scan,
OR-to-\texttt{ANY}, NOT NULL simplification, \dots). Writing grammar for each by
hand is slow, and the newest code is where regressions are. An LLM can read a
commit (message, diff, the regression test it adds) and write grammar for it.

Two designs we rejected:
\begin{itemize}
  \item \textbf{LLM as the generator} (Fuzz4All style): one model call per
    test case is 1{,}000$\times$ slower than our 17--23 test cases per second,
    and the output is unchecked text with no derivation tree (so no oracles, no
    tree reduction).
  \item \textbf{LLM as the oracle} (``is this result correct?''): LLMs cannot
    evaluate SQL on our data reliably; every hallucinated ``bug'' costs triage
    time.
\end{itemize}
Our design keeps the model \emph{off the hot path}: a few calls per commit
produce a grammar fragment; the grammar then generates thousands of
variations, each a typed tree the oracles can use. And nothing the model writes
is trusted: a gate on the real server decides.

\subsection{The pipeline}

\begin{enumerate}
  \item \textbf{Mine} (\file{src/pgfuzz/llm/commits.py}):
    \texttt{pgfuzz llm mine <pg clone> --since 2026-06-29 --until 26bd08fb
    --exclude bench/reinjected/commits.txt}. Keeps commits that touch
    query-processing code (optimizer, executor, parser, \file{utils/adt},
    rewrite, nodes, \dots) and skips typo/comment/pgindent/revert/warning
    commits. For each it stores a \texttt{CommitDigest}: sha, date, subject,
    body, files, changed function names, the SQL lines it added to the
    regression tests, and the \textbf{diff-coverage targets} (changed line
    numbers of every \texttt{.c} file).
  \item \textbf{Prompt} (\file{prompt.py}). A system prompt that is the same for
    every commit and therefore \textbf{cached} ($\approx$9{,}000 tokens): the
    \texttt{.pgg} format, every builtin with one line of documentation, the
    conventions (parenthesise operators, qualify columns, order-sensitive
    aggregates via \texttt{@sorted\_agg}), the schema, the type list, and the
    whole current grammar. The user message is only the commit digest.
  \item \textbf{Answer} in a strict JSON schema (structured output):
    \texttt{rationale} (which code paths changed and which SQL reaches them),
    \texttt{grammar\_fragment} (new alternatives as variants of existing rules,
    plus new rules), 1--5 \texttt{seeds}, \texttt{features}, and 0--3
    \texttt{rewrites} (equivalence templates for the metamorphic oracle).
  \item \textbf{Gate} (\file{gate.py}, below).
  \item \textbf{Repair}: a rejected proposal goes back to the model with the
    gate's findings -- the parse error, the top server error messages with
    counts, which alternatives were never derived, which changed lines were not
    covered -- up to 3 rounds.
  \item \textbf{Ledger}: each round is one line of \file{ledger.jsonl} (commit,
    round, tokens in/out/cached, cost, full gate report). Resumable; the
    source of RQ4 (\texttt{pgfuzz llm report}).
  \item \textbf{Use}: accepted fragments go to
    \file{results/llm/fragments/}, seeds to \file{results/llm/seeds/}, rewrites
    into the rewriter. The LLM arms of RQ1 are the grammar arm plus these files.
\end{enumerate}

\paragraph{The client} (\file{client.py}) has two modes. \texttt{api} calls
Claude through the Anthropic SDK. \texttt{manual} writes the prompt to
\file{prompts/<key>.md} and stops; a team member pastes it into Claude, saves
the JSON answer to \file{responses/<key>.json} and runs again -- no API key
needed. In both modes an existing response file is \emph{replayed}: reruns are
free and bit-for-bit reproducible, so the paper's numbers can be regenerated.
Token usage is stored with every response.

\subsection{The validation gate}
\label{sec:gate}

A fragment for one commit is accepted only if:
\begin{enumerate}
  \item \textbf{parse}: it parses together with the grammar, declares no new
    types and uses only known builtins;
  \item \textbf{reachability}: when we sample 150 test cases with the
    fragment's alternatives boosted (weight $\times 25$), they are derived in at
    least 5 statements (an alternative the engine can never complete is
    useless);
  \item \textbf{validity}: at least 60\% of the statements that use them run
    without error \emph{on the server} -- \pg decides, not the model;
  \item \textbf{seeds}: seed test cases run; failing seeds are dropped;
  \item \textbf{diff coverage} (on a gcov build): the fragment executes more of
    the commit's changed lines than the grammar without it, on the same sample
    size. Without a coverage build this is ``not measured'' and does not block.
\end{enumerate}
The gate is also a small fuzzing campaign aimed at the commit: bug candidates
hit while gating are recorded.

\paragraph{Diff coverage} (\file{diffcov.py}) $=$ changed lines executed /
changed lines that are executable, from gcov counters. Each backend writes
\file{.gcda} files when it exits, so the client reconnects before measuring.
If the coverage build is at a later revision than the commit, the changed line
numbers are mapped forward through \texttt{git diff}. All our builds are at the
same commit (\texttt{26bd08fb}), so the mapping is the identity.

\subsection{The rewrite gate (LLM-proposed metamorphic rewrites)}
\label{sec:rwgate}

A rewrite is a template with \texttt{\{e\}} and a type, for example
\texttt{CASE WHEN \{e\} IS NULL THEN NULL ELSE \{e\} END} for \texttt{int}.
A wrong rewrite would make the metamorphic oracle report false positives
forever, so the bar is higher:
\begin{enumerate}
  \item \textbf{static}: known type, contains \texttt{\{e\}}, balanced
    parentheses, no volatile function;
  \item \textbf{value probe}: evaluate
    \texttt{SELECT (template(v)) IS NOT DISTINCT FROM (v)} on the server for
    \texttt{v} = NULL, the type's edge values (min/max int, $\pm$Infinity,
    NaN, $-0$, empty string/array, \dots) and up to 6 data quantiles. One
    difference or error rejects it with the counterexample;
  \item \textbf{live checks}: apply \emph{only} this rewrite, with
    confirmation, to the expressions of 150 grammar test cases; accept only if
    it was applied at least 5 times, at most 20\% of the rewritten queries
    errored, and \textbf{no} confirmed mismatch.
\end{enumerate}

\begin{keybox}[Measured (rewrite gate on master)]
\begin{tabular}{lp{0.37\linewidth}l}
\toprule
Rewrite & Template & Result \\
\midrule
\texttt{null\_case} & \texttt{CASE WHEN \{e\} IS NULL THEN NULL ELSE \{e\} END} & accepted, 90/90 equal \\
\texttt{substr\_one} & \texttt{substr(\{e\}, 1)} (text) & accepted, 88/88 equal \\
\texttt{plus\_one} & \texttt{\{e\} + 1} (deliberately false) & rejected: 12 mismatches \\
\texttt{is\_true} & \texttt{\{e\} IS TRUE} (bool) & rejected by the \textbf{value probe} at NULL \\
\bottomrule
\end{tabular}

\texttt{e IS TRUE} is false, not NULL, when \texttt{e} is NULL. It passed 109
sampled live checks, because most predicates sit in \texttt{WHERE}, where NULL
and false filter the same rows. Sampling is not proof; the probe was added
because of this.
\end{keybox}

A confirmed mismatch is not thrown away: the gate stores the example with the
note ``not an identity, or a bug in \pg: triage the example''. (BUG-02 is
exactly a case where an identity, \texttt{e} $\equiv$ \texttt{(SELECT e)}, was
broken by \pg.)

\subsection{Leakage and cost}

\begin{itemize}
  \item \textbf{No leakage into RQ3.} The miner excludes the 23 commits whose
    fixes we revert, so the model never sees the fix of a bug it is later
    scored on. The mining window ends at the benchmark base \texttt{26bd08fb}.
    (Training-data leakage of older public commits cannot be excluded; it is
    a threat to validity, discussed in the paper.)
  \item \textbf{Cost.} Prices in \file{client.py}; the cached system prompt is
    read at a tenth of the input price. One round is about \$0.10; 30 commits
    with repair rounds about \$5--10 (estimate, to be replaced by the ledger
    total).
  \item \textbf{Status.} The pipeline is tested end to end with a live gate
    (it accepted a valid fragment, 251/251 statements valid, and rejected an
    invalid one with the server's own error message). The real extension run
    over $\approx$30 commits is Phase~1 of the execution plan; the RQ3 smoke
    ``+LLM'' arm ran \emph{without} fragments, so its result equals
    Grammar+Oracles, as expected.
\end{itemize}

\subsection{RQ4: what we will report}
Accepted / rejected fragments per gate check; repair rounds needed; diff
coverage with vs without the fragment, per commit; tokens and dollars; the
ablation of RQ1 (\texttt{llm\_frag}: fragments only, \texttt{llm\_seeds}: seeds
only, \texttt{grammar\_llm}: both) against \texttt{grammar\_mm}; bugs whose
test case used an LLM alternative (the tree records it).
